\begin{problem}
    En la parte superior de un poste de \SI{3}{m} de altura perpendicular al suelo, se atan dos cuerdas de manera tensa, quedando de \SI{4}{m} cada una, tal como se representa en la figura adjunta.

    \begin{center}
    (imagen)
    \end{center}

    Para determinar la distancia $x$ con respecto al poste, se realiza el siguiente procedimiento, cometiéndose un error.

    \begin{itemize}
        \item Paso 1: para calcular la longitud $x$ se decide utilizar el teorema de Pitágoras, por lo que se reemplazan los valores, obteniéndose $x^2 + 3^2 = 4^2$.
        \item Paso 2: se calculan las potencias en la ecuación anterior, obteniéndose $x^2 + 9 = 16$.
        \item Paso 3: se despeja $x$, obteniéndose $x = \sqrt{16 + 9}$.
        \item Paso 4: se calcula la raíz cuadrada del paso anterior, obteniéndose la distancia de \SI{5}{m}.
    \end{itemize}

    ¿En cuál de los pasos se cometió el error?

    \begin{enumerate}[label=\Alph*.]
        \item En el Paso 1
        \item En el Paso 2
        \item En el Paso 3
        \item En el Paso 4
    \end{enumerate}
\end{problem}
